\begin{lemma}
\label{lemma-morphism-ringed-etale-topoi-affines}
Let $X$, $Y$ be affine schemes.
Let
$$
(g, g^\#) :
(\Sh(X_\etale), \mathcal{O}_X)
\longrightarrow
(\Sh(Y_\etale), \mathcal{O}_Y)
$$
be a morphism of locally ringed topoi. Then there exists a
unique morphism of schemes $f : X \to Y$ such that
$(g, g^\#)$ is $2$-isomorphic to $(f_{small}, f_{small}^\sharp)$,
see
Modules on Sites,
Definition \ref{sites-modules-definition-2-morphism-ringed-topoi}.
\end{lemma}

\begin{proof}
In this proof we write $\mathcal{O}_X$ for the structure sheaf
of the small \'etale site $X_\etale$, and similarly for
$\mathcal{O}_Y$. Say $Y = \Spec(B)$ and $X = \Spec(A)$. Since
$B = \Gamma(Y_\etale, \mathcal{O}_Y)$,
$A = \Gamma(X_\etale, \mathcal{O}_X)$
we see that $g^\sharp$ induces a ring map $\varphi : B \to A$.
Let $f = \Spec(\varphi) : X \to Y$ be the corresponding morphism
of affine schemes. We will show this $f$ does the job.

\medskip\noindent
Let $V \to Y$ be an affine scheme \'etale over $Y$. Thus we may write
$V = \Spec(C)$ with $C$ an \'etale $B$-algebra. We can write
$$
C = B[x_1, \ldots, x_n]/(P_1, \ldots, P_n)
$$
with $P_i$ polynomials such that $\Delta = \det(\partial P_i/ \partial x_j)$
is invertible in $C$, see for example
Algebra, Lemma \ref{algebra-lemma-etale-standard-smooth}.
If $T$ is a scheme over $Y$, then a $T$-valued point of $V$ is given by
$n$ sections of $\Gamma(T, \mathcal{O}_T)$ which satisfy the polynomial
equations $P_1 = 0, \ldots, P_n = 0$. In other words, the sheaf $h_V$
on $Y_\etale$ is the equalizer of the two maps
$$
\xymatrix{
\prod\nolimits_{i = 1, \ldots, n} \mathcal{O}_Y
\ar@<1ex>[r]^a \ar@<-1ex>[r]_b &
\prod\nolimits_{j = 1, \ldots, n} \mathcal{O}_Y
}
$$
where $b(h_1, \ldots, h_n) = 0$ and
$a(h_1, \ldots, h_n) =
(P_1(h_1, \ldots, h_n), \ldots, P_n(h_1, \ldots, h_n))$.
Since $g^{-1}$ is exact we conclude that the top row of the
following solid commutative diagram is an equalizer diagram as well:
$$
\xymatrix{
g^{-1}h_V \ar[r] \ar@{..>}[d] &
\prod\nolimits_{i = 1, \ldots, n} g^{-1}\mathcal{O}_Y
\ar@<1ex>[r]^{g^{-1}a} \ar@<-1ex>[r]_{g^{-1}b} \ar[d]^{\prod g^\sharp} &
\prod\nolimits_{j = 1, \ldots, n} g^{-1}\mathcal{O}_Y \ar[d]^{\prod g^\sharp}\\
h_{X \times_Y V} \ar[r] &
\prod\nolimits_{i = 1, \ldots, n} \mathcal{O}_X
\ar@<1ex>[r]^{a'} \ar@<-1ex>[r]_{b'} &
\prod\nolimits_{j = 1, \ldots, n} \mathcal{O}_X  \\
}
$$
Here $b'$ is the zero map and $a'$ is the map defined by the
images $P'_i = \varphi(P_i) \in A[x_1, \ldots, x_n]$ via the same
rule
$a'(h_1, \ldots, h_n) =
(P'_1(h_1, \ldots, h_n), \ldots, P'_n(h_1, \ldots, h_n))$.
that $a$ was defined by. The commutativity of the diagram follows from
the fact that $\varphi = g^\sharp$ on global sections. The lower
row is an equalizer diagram also, by exactly the same arguments as
before since $X \times_Y V$ is the affine scheme
$\Spec(A \otimes_B C)$ and
$A \otimes_B C = A[x_1, \ldots, x_n]/(P'_1, \ldots, P'_n)$.
Thus we obtain a unique dotted arrow
$g^{-1}h_V \to h_{X \times_Y V}$ fitting into the diagram

\medskip\noindent
We claim that the map of sheaves $g^{-1}h_V \to h_{X \times_Y V}$
is an isomorphism. Since the small \'etale site of $X$ has enough points
(Theorem \ref{theorem-exactness-stalks})
it suffices to prove this on stalks. Hence let $\overline{x}$ be a
geometric point of $X$, and denote $p$ the associate point of the
small \'etale topos of $X$. Set $q = g \circ p$. This is a point of
the small \'etale topos of $Y$. By
Lemma \ref{lemma-points-small-etale-site}
we see that $q$ corresponds to a geometric point $\overline{y}$ of
$Y$. Consider the map of stalks
$$
(g^\sharp)_p :
(\mathcal{O}_Y)_{\overline{y}} =
\mathcal{O}_{Y, q} =
(g^{-1}\mathcal{O}_Y)_p
\longrightarrow
\mathcal{O}_{X, p} =
(\mathcal{O}_X)_{\overline{x}}
$$
Since $(g, g^\sharp)$ is a morphism of {\it locally} ringed topoi
$(g^\sharp)_p$ is a local ring homomorphism of strictly henselian
local rings. Applying localization to the big commutative diagram above and
Algebra, Lemma \ref{algebra-lemma-strictly-henselian-solutions}
we conclude that $(g^{-1}h_V)_p \to (h_{X \times_Y V})_p$ is an isomorphism
as desired.

\medskip\noindent
We claim that the isomorphisms $g^{-1}h_V \to h_{X \times_Y V}$ are
functorial. Namely, suppose that $V_1 \to V_2$ is a morphism of affine
schemes \'etale over $Y$. Write
$V_i = \Spec(C_i)$ with
$$
C_i = B[x_{i, 1}, \ldots, x_{i, n_i}]/(P_{i, 1}, \ldots, P_{i, n_i})
$$
The morphism $V_1 \to V_2$ is given by a $B$-algebra map $C_2 \to C_1$
which in turn is given by some polynomials
$Q_j \in B[x_{1, 1}, \ldots, x_{1, n_1}]$ for $j = 1, \ldots, n_2$.
Then it is an easy matter to show that the diagram of sheaves
$$
\xymatrix{
h_{V_1} \ar[d] \ar[r] & \prod_{i = 1, \ldots, n_1} \mathcal{O}_Y
\ar[d]^{Q_1, \ldots, Q_{n_2}}\\
h_{V_2} \ar[r] & \prod_{i = 1, \ldots, n_2} \mathcal{O}_Y
}
$$
is commutative, and pulling back to $X_\etale$ we obtain the
solid commutative diagram
$$
\xymatrix{
g^{-1}h_{V_1} \ar@{..>}[dd] \ar[rrd] \ar[r] &
\prod_{i = 1, \ldots, n_1} g^{-1}\mathcal{O}_Y
\ar[dd]^{g^\sharp}
\ar[rrd]^{Q_1, \ldots, Q_{n_2}} \\
& & g^{-1}h_{V_2} \ar@{..>}[dd] \ar[r] &
\prod_{i = 1, \ldots, n_2} g^{-1}\mathcal{O}_Y
\ar[dd]^{g^\sharp} \\
h_{X \times_Y V_1} \ar[r] \ar[rrd] &
\prod\nolimits_{i = 1, \ldots, n_1} \mathcal{O}_X
\ar[rrd]^{Q'_1, \ldots, Q'_{n_2}} \\
& & h_{X \times_Y V_2} \ar[r] &
\prod\nolimits_{i = 1, \ldots, n_2} \mathcal{O}_X
}
$$
where $Q'_j \in A[x_{1, 1}, \ldots, x_{1, n_1}]$ is the image of
$Q_j$ via $\varphi$. Since the dotted arrows exist, make the
two squares commute, and the horizontal arrows are injective
we see that the whole diagram commutes. This proves functoriality
(and also that the construction of $g^{-1}h_V \to h_{X \times_Y V}$
is independent of the choice of the presentation, although we
strictly speaking do not need to show this).

\medskip\noindent
At this point we are able to show that $f_{small, *} \cong g_*$.
Namely, let $\mathcal{F}$ be a sheaf on $X_\etale$. For every
$V \in \Ob(X_\etale)$ affine we have
\begin{align*}
(g_*\mathcal{F})(V)
& =
\Mor_{\Sh(Y_\etale)}(h_V, g_*\mathcal{F}) \\
& =
\Mor_{\Sh(X_\etale)}(g^{-1}h_V, \mathcal{F}) \\
& =
\Mor_{\Sh(X_\etale)}(h_{X \times_Y V}, \mathcal{F}) \\
& =
\mathcal{F}(X \times_Y V) \\
& =
f_{small, *}\mathcal{F}(V)
\end{align*}
where in the third equality we use the isomorphism
$g^{-1}h_V \cong h_{X \times_Y V}$ constructed above. These isomorphisms
are clearly functorial in $\mathcal{F}$ and functorial in $V$
as the isomorphisms $g^{-1}h_V \cong h_{X \times_Y V}$ are functorial.
Now any sheaf on $Y_\etale$ is determined by the restriction
to the subcategory of affine schemes
(Topologies, Lemma \ref{topologies-lemma-alternative}),
and hence we obtain an isomorphism of functors $f_{small, *} \cong g_*$
as desired.

\medskip\noindent
Finally, we have to check that, via the isomorphism
$f_{small, *} \cong g_*$ above, the maps $f_{small}^\sharp$ and
$g^\sharp$ agree. By construction this is already the case for the
global sections of $\mathcal{O}_Y$, i.e., for the elements of $B$.
We only need to check the result on
sections over an affine $V$ \'etale over $Y$ (by
Topologies, Lemma \ref{topologies-lemma-alternative}
again). Writing
$V = \Spec(C)$, $C = B[x_i]/(P_j)$ as before it suffices
to check that the coordinate functions $x_i$ are mapped to
the same sections of $\mathcal{O}_X$ over $X \times_Y V$.
And this is exactly what it means that the diagram
$$
\xymatrix{
g^{-1}h_V \ar[r] \ar@{..>}[d] &
\prod\nolimits_{i = 1, \ldots, n} g^{-1}\mathcal{O}_Y
\ar[d]^{\prod g^\sharp} \\
h_{X \times_Y V} \ar[r] &
\prod\nolimits_{i = 1, \ldots, n} \mathcal{O}_X
}
$$
commutes. Thus the lemma is proved.
\end{proof}